\documentclass[DIV=12,%
               BCOR=0mm,%
               fontsize=10pt,%
               oneside,%
               paper=210mm:11in]{scrbook}
\usepackage{fontspec}
\usepackage{polyglossia}
\setmainfont[Mapping=tex-text]{Linux Libertine O}
% these are not used but prevents XeTeX to barf
\setsansfont[Mapping=tex-text,Scale=MatchLowercase]{DejaVu Sans}
\setmonofont[Mapping=tex-text,Scale=MatchLowercase]{DejaVu Sans Mono}

\setmainlanguage{english}





\usepackage{microtype} % you need an *updated* texlive 2012, but harmless
\usepackage{graphicx}
\usepackage{alltt}
\usepackage{verbatim}
% http://tex.stackexchange.com/questions/3033/forcing-linebreaks-in-url
\PassOptionsToPackage{hyphens}{url}\usepackage[hyperfootnotes=false,hidelinks,breaklinks=true]{hyperref}
\usepackage{bookmark}
\usepackage[stable]{footmisc}
\usepackage{enumerate}
\usepackage{tabularx}
\usepackage[normalem]{ulem}
\usepackage{wrapfig}
\usepackage{indentfirst}
% remove the numbering
\setcounter{secnumdepth}{-2}

% remove labels from the captions
\renewcommand*{\captionformat}{}
\renewcommand*{\figureformat}{}
\renewcommand*{\tableformat}{}
\KOMAoption{captions}{belowfigure,nooneline}
\addtokomafont{caption}{\centering}

% avoid breakage on multiple <br><br> and avoid the next [] to be eaten
\newcommand*{\forcelinebreak}{\strut\\{}}

\newcommand*{\hairline}{%
  \bigskip%
  \noindent \hrulefill%
  \bigskip%
}

% reverse indentation for biblio and play

\newenvironment*{amusebiblio}{
  \leftskip=\parindent
  \parindent=-\parindent
  \smallskip
  \indent
}{\smallskip}

\newenvironment*{amuseplay}{
  \leftskip=\parindent
  \parindent=-\parindent
  \smallskip
  \indent
}{\smallskip}

\newcommand*{\Slash}{\slash\hspace{0pt}}

% global style
\pagestyle{plain}
\addtokomafont{disposition}{\rmfamily}
% forbid widows/orphans
\clubpenalty=10000
\widowpenalty=10000
\frenchspacing
\sloppy

\title{The Text::Amuse markup manual}
\date{2012}
\author{}
\subtitle{The writer's guide}

\begin{document}


\maketitle




\tableofcontents
% start a new right-handed page
\cleardoublepage




Copyright (C) 2004, 2005, 2006, 2007, 2008, 2009, 2010 Free Software
Foundation, Inc.



\begin{quote}


Permission is granted to copy, distribute and\Slash{}or modify this
document under the terms of the GNU Free Documentation License,
Version 1.2 or any later version published by the Free Software
Foundation; with no Invariant Sections, no Front-Cover Texts, and
no Back-Cover Texts.  A copy of the license is included in the
section entitled “GNU Free Documentation License.”



\end{quote}

\part{The Muse markup}


A Muse document uses special, contextual markup rules to determine how
to format the output result.  For example, if a paragraph is indented,
Muse assumes it should be quoted. Indentation is heavily used to
determine if the paragraph is something different from “normal” text.


There are not too many markup rules, and all of them strive to be as
simple as possible so that you can focus on document creation, rather
than formatting.


This document describes Muse, which was written by John
Wiegley, then maintained by Michael Olson and edited for this
implementation by Marco Pessotto.

\chapter{Paragraphs}


Paragraphs in Muse must be separated by a blank line.

\section{Centered and right aligned paragraphs and quotations}


A line that begins with six or more columns of whitespace
(but less than 20) indicates a centered paragraph.


Alternatively, you can use the \texttt{<center>} tag to surround regions that
are to be published as centered paragraphs.



\begin{center}

Like this


\end{center}


But if the line begins with more than 20 spaces, you'll have a right
aligned paragraph. This is handy for signatures.



\begin{flushright}

Like this right-aligned one.


\end{flushright}


The same result is done with the \texttt{right} tag.


Using the tags:



\begin{center}

Like this centered one


\end{center}



\begin{flushright}

And this is right


\end{flushright}


This is a quotation:



\begin{quote}


But if a line begins with whitespace (at least 2 spaces), though
less than six columns, it indicates a quoted paragraph.
Alternatively, you can use the \texttt{<quote>} tag to surround regions that
are to be published as quoted paragraphs.



\end{quote}

\subsection{Source}


\begin{alltt}
A line that begins with six or more columns of whitespace
(but less than 20) indicates a centered paragraph.

Alternatively, you can use the =<center>= tag to surround regions that
are to be published as centered paragraphs.

           Like this

But if the line begins with more than 20 spaces, you'll have a right
aligned paragraph. This is handy for signatures.

                                     Like this right-aligned one.

The same result is done with the =right= tag.

Using the tags:

<center>
Like this centered one
</center>

<right>
And this is right
</right>

This is a quotation:

  But if a line begins with whitespace (at least 2 spaces), though
  less than six columns, it indicates a quoted paragraph.
  Alternatively, you can use the =<quote>= tag to surround regions that
  are to be published as quoted paragraphs.

\end{alltt}


\textbf{Please keep in mind that indentation must be consistent if you prefer
to keep the lines short} and break them inserting a new line. Long
lines (using the rule 1 line, 1 paragraph) are perfectly fine. The
rule of paragraphs separated by blank lines still apply, though.


Also, \texttt{<tags>} which start and stop blocks, \textbf{must be placed on a line by
themselves}.

\section{Literal paragraphs}


The \texttt{<example>} tag is used for examples, where whitespace should be
preserved, the text rendered in monospace, and any characters special
to the output style escaped.


Example:


\begin{alltt}
 <example>
 The =<example>= tag is used for examples, where whitespace
 should be preserved, the text rendered in monospace, and any
 characters special to the output style escaped.
 </example>

\end{alltt}


There is no \texttt{<literal>} tag as in the original Muse markup, because it's
not a private tool and will be exposed to the internet.

\section{Line breaks}


If you need a line break, then use the \texttt{br} tag.  Most of the time
this tag is unnecessary, because Muse will automatically detect
paragraphs by means of blank lines.  If you want to preserve newlines in
several lines of text, then use \texttt{verse} markup instead.


\begin{alltt}
This line will break <br> And continue

\end{alltt}


Yields:


This line will break \forcelinebreak  And continue.



\bigskip


If you want to add a blank, line, put the \texttt{br} tag on a line by itself:


\begin{alltt}
Here we add a blank line

<br>

Here we go.

\end{alltt}


Here we add a blank line



\bigskip


Here we go.

\section{Page breaks}


If you put exactly five “*” separated by whitespace on a line by
itself, indented by 6 or more spaces (like a centered paragraph),
you'll get a page break in the PDF.


This code will break the page.


\begin{alltt}
       * * * * *

\end{alltt}



\begin{center}

* * * * *


\cleardoublepage


\end{center}


Anyway, using three of them is just a decorator and it's not threated specially.



\begin{center}

* * *


\end{center}


It's just a centered paragraph with 3 “*”.

\chapter{Levels of headings}


A heading becomes a chapter or section in printed output - depending on
the style.  To indicate a heading, start a new paragraph with one or
more asterices, followed by a space and the heading title.  Then begin
another paragraph to enter the text for that section.


All levels of headings will be published. There is support for 5 levels.


The first level is a "part", and should be used only for larger texts. In this document is used for the License and for the main title.


The second level is a "chapter". It starts a new page on the PDF output.


The third level is undoubtely the most used. It usually separate a
section of an article. For example the "Literal paragraph" above.


The forth level goes down further.

\subsection{Forth level, a “subsection”}


The fifth level is very, very low. It's handy to create description
lists, too. It's not dumped on the Table of Contents.

\subsubsection{item 1}


A description

\subsubsection{item 2}


Another description

\subsection{The example of levels}


\begin{alltt}
* First level, aka part

** Second level, aka chapter

*** Third level, aka section

**** Fourth level, aka subsection

***** Fifth level, aka subsubsection

\end{alltt}

\chapter{Directives at the beginning of a document}


Directives are lines beginning with the `\#' character that come before
any paragraphs or sections in the document.


Directives are of the form \texttt{\#directive content of directive}.


You can use any combination of uppercase and lowercase letters for
directives, even if the directive is not in the list below. The
directives are completely arbitrary. You can put there whatever you
want. It's the template job to pick them up. In the templates shipped
with this bunlde there is support for the following directives:



\begin{quote}


The following is a list of directives that Muse uses.



\end{quote}


\begin{itemize}
\item\relax 
\texttt{\#author}


The author of the text.



\item\relax 
\texttt{\#title}


The title of the document



\item\relax 
\texttt{\#lang}


The language code of the document (2 or 3 letters). Defaults to \texttt{en}.



\item\relax 
\texttt{\#LISTtitle}


This directive is used (defaulting to \texttt{\#title}) to alphabetically
sort the titles. It's handy if you want, for example, sort "A
title" under "T" and not under "A".


In this case you may write \texttt{\#LISTtitle Title}



\item\relax 
\texttt{\#subtitle}


The subtitle (if any)



\item\relax 
\texttt{\#SORTauthors}


If not provided, this default to \texttt{\#author}. It's a list
separated by semicolons or commas with the various authors. While
\texttt{\#author} affects the display only, this one is used to index the
document.



\item\relax 
\texttt{\#SORTtopics}


As for authors, it's a list (comma- or
semicolon-separated) list of topics for the current text. Used to
index the document.



\item\relax 
\texttt{\#date}


The \textbf{year} of publishing of the document. More information
should be provided in the \texttt{\#notes} directive.



\item\relax 
\texttt{\#notes}


This directive is used for additional information here (original
title, translators, credits, etc).



\item\relax 
\texttt{\#source}


This directive is used for the source or the text (url, scanned from
original, original contribution, etc.). The preferred format is
“Retrieved on March 8, 2012 from \href{http://url.org}{the url}”




\end{itemize}
\chapter{Bold, italicized and emphatized text}


To emphasize text, surround it with certain specially recognized
characters.


The following example will produce:



\begin{quote}


\emph{emphasis}
\textbf{strong emphasis}
\textbf{\emph{very strong emphasis}}
\texttt{verbatim and monospace}



\end{quote}


\begin{alltt}
     *emphasis*
     **strong emphasis**
     ***very strong emphasis***
     =verbatim and monospace=

\end{alltt}


Each of these forms may span
multiple lines, but not multiple paragraphs.


You can also use the \texttt{<code>} tag to indicate verbatim and monospace
text. This is handy for regions that have an "=" in them.


If the “*” confuse you or the preview is screwed up, you can use
inline tag \texttt{<em>} and \texttt{<strong>}, which are guarantee to work in any case.


The above example rewritten with tags:


\begin{alltt}
     <em>emphasis</em>
     <strong>strong emphasis</strong>
     <strong><em>very strong emphasis</em></strong>
     <code>verbatim and monospace</code>

\end{alltt}


And produces the same thing.



\begin{quote}


\emph{emphasis}
\textbf{strong emphasis}
\textbf{\emph{very strong emphasis}}
\texttt{verbatim and monospace}



\end{quote}


Please note that there is no support for the underline. Underlining is
an handwritten substitute for the italics. You simply don't need it.


Also, small caps are missing, mainly because on the HTML they look
awful and a very few fonts have decent small caps.


Other tags are \texttt{<sub>} and \texttt{<sup>} for subscript and superscript. And
there is also a \texttt{<del>} tag for overstriking.


\begin{alltt}
This is the <sup>superscript</sup> and this is a
<sub>subscript</sub>, and this is <del>something deleted and
overstriked</del>

\end{alltt}


This is the \textsuperscript{superscript} and this is a
\textsubscript{subscript}, and this is \sout{something deleted and
overstriked}

\section{Making notes to be shown at the end}


A footnote reference is simply a number in square brackets.  To define
the footnote, place this definition starting the line with a digit in
square brackets.


\begin{alltt}
This is the text, and we refer to a footnote [1]

Here the text continues.

[1] This footnote
    spans more lines in the source

 [2] But this is not, because of the initial whitespace.

[3] Footnotes which don't have a referrer will disappear on the PDF
    output and preserverd in the HTML. But will lead to incorrect
    code, as it will point to a non-existent anchor

\end{alltt}


This is the result:


\hairline


This is the text, and we refer to a footnote\footnote{This footnote spans more lines in the source}


Here the text continues.


 [2] But this is not, because of the initial whitespace.


\hairline


You can break the footnotes lines (even if it's not reccomended), but keep the indentation consistent, as shown above

\section{Indicating poetic stanzas}


Poetry requires that whitespace be preserved, but without resorting to
monospace.  To indicate this, use the following markup, reminiscent of
email quotations, or use the \texttt{verse} tag.


\begin{alltt}
> A line of Emacs verse;
>   forgive its being so terse.

<verse>
     A line of Emacs verse;
       forgive its being so terse.
</verse>

\end{alltt}


This yields:



\begin{verse}
A line of Emacs verse;\forcelinebreak
~~forgive its being so terse.
\end{verse}



\begin{verse}
~~~~~A line of Emacs verse;\forcelinebreak
~~~~~~~forgive its being so terse.
\end{verse}


Multiple stanzas may be included in one set of <verse> tags, as
follows.


\begin{alltt}
<verse>
A line of Emacs verse;
  forgive its being so terse.

In terms of terse verse,
  you could do worse.
</verse>

Or this

> A line of Emacs verse;
>   forgive its being so terse.
>
> In terms of terse verse,
>   you could do worse.


\end{alltt}



\begin{verse}
A line of Emacs verse;\forcelinebreak
~~forgive its being so terse.

In terms of terse verse,\forcelinebreak
~~you could do worse.
\end{verse}


Or this



\begin{verse}
A line of Emacs verse;\forcelinebreak
~~forgive its being so terse.

In terms of terse verse,\forcelinebreak
~~you could do worse.
\end{verse}

\chapter{Lists}


Lists are given using special characters at the beginning of a line.
Whitespace must occur before bullets or numbered items, to distinguish
from the possibility of those characters occurring in a real sentence.


\begin{alltt}
Normal text.

      - bullet item one
      - bullet item two

An enumerated list follows.

      1. Enum item one
      2. Enum item two

A list with roman numbering

 i.   First
 ii.  Second
 iii. Third

A list with upper roman numbering

 I.   First
 II.  Second
 III. Third

A list with upper letters

 A. first
 B. second
 C. third

A list with lower letters

 a. first
 b. second
 c. third

\end{alltt}


Please note the consistent indentation, \textbf{expecially} for roman numbering.


Normal text.


\begin{itemize}
\item\relax 
bullet item one



\item\relax 
bullet item two




\end{itemize}

An enumerated list follows.


\begin{enumerate}[1.]
\item\relax 
Enum item one



\item\relax 
Enum item two




\end{enumerate}

A list with roman numbering


\begin{enumerate}[i.]
\item\relax 
First



\item\relax 
Second



\item\relax 
Third




\end{enumerate}

A list with upper roman numbering


\begin{enumerate}[I.]
\item\relax 
First



\item\relax 
Second



\item\relax 
Third




\end{enumerate}

A list with upper letters


\begin{enumerate}[A.]
\item\relax 
first



\item\relax 
second



\item\relax 
third




\end{enumerate}

A list with lower letters


\begin{enumerate}[a.]
\item\relax 
first



\item\relax 
second



\item\relax 
third




\end{enumerate}

While the original Emacs Muse had the description lists, this
implementation has not. You should use the fifth level of headings for
those. The outcome is mostly the same (yes, you can't nest them in
lists, but crazy nesting is beyond the scope of this markup).

\section{Nested lists}


It is possible to nest lists of the same or different kinds.  The
"level" of the list is determined by the amount of initial whitespace.


\begin{alltt}
Normal text.

 - Level 1, bullet item one
   1. Level 2, enum item one
   2. Level 2, enum item two
 - Level 1, bullet item two
   1. Level 2, enum item one
   2. Level 2, enum item two
      i.  Level 3, enum item i
      ii. Level 3, enum item ii
   3. Level 2, enum item three
 - Back to Level 1, third bullet
   a. Level 2, enum item “a”
   b. Level 2, enum item “b”
      I. Level 3, enum item “I”
 - Back to the bullets

\end{alltt}


Normal text.


\begin{itemize}
\item\relax 
Level 1, bullet item one


\begin{enumerate}[1.]
\item\relax 
Level 2, enum item one



\item\relax 
Level 2, enum item two




\end{enumerate}


\item\relax 
Level 1, bullet item two


\begin{enumerate}[1.]
\item\relax 
Level 2, enum item one



\item\relax 
Level 2, enum item two


\begin{enumerate}[i.]
\item\relax 
Level 3, enum item i



\item\relax 
Level 3, enum item ii




\end{enumerate}


\item\relax 
Level 2, enum item three




\end{enumerate}


\item\relax 
Back to Level 1, third bullet


\begin{enumerate}[a.]
\item\relax 
Level 2, enum item “a”



\item\relax 
Level 2, enum item “b”


\begin{enumerate}[I.]
\item\relax 
Level 3, enum item “I”




\end{enumerate}



\end{enumerate}


\item\relax 
Back to the bullets




\end{itemize}
\section{Breaking list items}


If you want to break up a line within any list type, just put one blank
line between the end of the previous line and the beginning of the next
line, using the same amount of initial indentation.


Keep in mind that if you put random indentation you'll get random and
probably unexpected results (but it should not crash -- if it does,
please contact me).


Also, you can be lazy with numbered list. The parser actually doesn't
care if you number them properly, or just do something like that.


\begin{alltt}
 1. first
 1. second
 1. third

or

 a. first
 a. second
 a. third

\end{alltt}


There results will always be:


\begin{enumerate}[1.]
\item\relax 
first



\item\relax 
second



\item\relax 
third




\end{enumerate}

or


\begin{enumerate}[a.]
\item\relax 
first



\item\relax 
second



\item\relax 
third




\end{enumerate}

If you need to switch type while the list is running, put a \texttt{br} tag
between them, or some invisible material like a non-breaking space.


\begin{alltt}
This won't get you what you expect:

 1. first
 a. first item

So you'll have to do:

 1. first

<br>

 a. first item

\end{alltt}


Which yields:


\hairline


This won't get you what you expect:


\begin{enumerate}[1.]
\item\relax 
first



\item\relax 
first item




\end{enumerate}

So you'll have to do:


\begin{enumerate}[1.]
\item\relax 
first




\end{enumerate}


\bigskip


\begin{enumerate}[a.]
\item\relax 
first item




\end{enumerate}

\hairline

\section{Complete example}


\begin{alltt}
Normal text.

 - Level 1, bullet item one, this is the first paragraph. I can break
   the line, keeping the same amount of indentation

   Here I have the same amount of indentation, and it continues the
   item above.

   1. Level 2, enum item one. i can break the line, keeping the same
      amount of indentation

      Here I have the same amount of indentation, and it continues the
      item above.

   2. Level 2, enum item two
      which continues

      Here I have the same amount of indentation, and it continues the
      item above.

 - Level 1, bullet item two
   which continues

   Here I have the same amount of indentation, and it continues the
   item above.

   1. Level 2, enum item one
      which continues

      Here I have the same amount of indentation, and it continues the
      item above.

   2. Level 2, enum item two
      which continues

      Here I have the same amount of indentation, and it continues the
      item above.

      i.  Level 3, enum item i

          Here I have the same amount of indentation, and it continues
          the item above.

      ii. Level 3, enum item ii

          Here I have the same amount of indentation, and it continues
          the item above.

   3. Level 2, enum item three
      which continues

      Here I have the same amount of indentation, and it continues the
      item above.

 - Back to Level 1, third bullet

   Here I have the same amount of indentation, and it continues the
   item above.

   a. Level 2, enum item “a”
      which continues

      Here I have the same amount of indentation, and it continues the
      item above.

   b. Level 2, enum item “b”
      which continues

      Here I have the same amount of indentation, and it continues the
      item above.

         I. Level 3, enum item “I”

            Here I have the same amount of indentation, and it continues the
            item above.

 - Back to the bullets

   Here I have the same amount of indentation, and it continues the
   item above.


\end{alltt}


Normal text.


\begin{itemize}
\item\relax 
Level 1, bullet item one, this is the first paragraph. I can break
the line, keeping the same amount of indentation


Here I have the same amount of indentation, and it continues the
item above.


\begin{enumerate}[1.]
\item\relax 
Level 2, enum item one. i can break the line, keeping the same
amount of indentation


Here I have the same amount of indentation, and it continues the
item above.



\item\relax 
Level 2, enum item two
which continues


Here I have the same amount of indentation, and it continues the
item above.




\end{enumerate}


\item\relax 
Level 1, bullet item two
which continues


Here I have the same amount of indentation, and it continues the
item above.


\begin{enumerate}[1.]
\item\relax 
Level 2, enum item one
which continues


Here I have the same amount of indentation, and it continues the
item above.



\item\relax 
Level 2, enum item two
which continues


Here I have the same amount of indentation, and it continues the
item above.


\begin{enumerate}[i.]
\item\relax 
Level 3, enum item i


Here I have the same amount of indentation, and it continues
the item above.



\item\relax 
Level 3, enum item ii


Here I have the same amount of indentation, and it continues
the item above.




\end{enumerate}


\item\relax 
Level 2, enum item three
which continues


Here I have the same amount of indentation, and it continues the
item above.




\end{enumerate}


\item\relax 
Back to Level 1, third bullet


Here I have the same amount of indentation, and it continues the
item above.


\begin{enumerate}[a.]
\item\relax 
Level 2, enum item “a”
which continues


Here I have the same amount of indentation, and it continues the
item above.



\item\relax 
Level 2, enum item “b”
which continues


Here I have the same amount of indentation, and it continues the
item above.


\begin{enumerate}[I.]
\item\relax 
Level 3, enum item “I”


Here I have the same amount of indentation, and it continues the
item above.




\end{enumerate}



\end{enumerate}


\item\relax 
Back to the bullets


Here I have the same amount of indentation, and it continues the
item above.




\end{itemize}
\chapter{Generation of data tables}


Only very simple tables are supported.  The syntax is as follows (just keep the indentation consistent and separate each cell by one or more textbars.


\begin{alltt}
     Triple bars \textbar{}\textbar{}\textbar{} Separate footer fields
     Double bars  \textbar{}\textbar{} Separate header fields
     Single bars   \textbar{} Separate body fields
     Here are more \textbar{} body fields
     \textbar{}+ This is the caption +\textbar{}

\end{alltt}

\begin{table}[htbp!]
 \begin{minipage}[t]{\textwidth}
\begin{tabularx}{\textwidth}{|X|X|}
\hline
\relax      Double bars   &  Separate header fields  \\
\hline
\relax      Single bars    &  Separate body fields  \\
\relax      Here are more  &  body fields  \\
\hline
\relax      Triple bars  &  Separate footer fields  \\
\hline
\end{tabularx}

\caption[]{This is the caption}
\end{minipage}
\end{table}


The ordering of the footer, header and table body blocks is irrelevant
for the output (HTML requires you put first the header, then the
footer, then the body). Ordering of the single rows is of course
preserved. From this we come to the conclusion that you can't normally
typeset the pipe symbol without triggering a list. In \texttt{<example>} you
can type it freely. Inside the cells you can do pretty much what you
want (beside headers and lists): mark them up freely.

\section{Floating tables}


When composing tables, please keep in mind the following:


\begin{itemize}
\item\relax 
tables are never splat over pages. If they don't fit a page, they
will run off the page. So keep them short or split them yourself.



\item\relax 
tables \emph{without} a caption are inserted at the point they are in the
source. Table \emph{with} a caption are converted into a float, so they
are in no way guaranteed they will appear in the same exact point
of the source. Anyway, the result is way better, and you have a
caption, so you can refer to them in the text.




\end{itemize}
\chapter{Hyperlinks with or without description (and images)}


A hyperlink can reference a URL or a place on the same document.  In
addition, descriptive text can be specified, which should be displayed
rather than the link text in output styles that supports link
descriptions.  The syntax is as follows.


\begin{alltt}
     [[link target][link description]]
     [[link target without description]]

So, the home of this project is
[[http://theanarchistlibrary.org][The Anarchist Library]],
which can be found at [[http://theanarchistlibrary.org]]
Bare links will not get the hyperlinking.
So http://thisisspam.org won't get the hyperlinking.

\end{alltt}


So, the home of this project is
\href{http://theanarchistlibrary.org}{The Anarchist Library},
which can be found at \url{http://theanarchistlibrary.org}
Bare links will not get the hyperlinking.
So http:\Slash{}\Slash{}thisisspam.org won't get the hyperlinking.

\section{Images}


Images are special case of this kind of linking.


\begin{alltt}
[[logo.png]]

\end{alltt}


We assume that we have the “logo.png” file on the same
directory of the file, and the result is:


\begin{figure}[htbp!]
\centering
\includegraphics[width=\textwidth]{logo.png}
\end{figure}


Now, let's add a caption.


\begin{alltt}
[[logo.png][This is *our* logo]]

\end{alltt}


\begin{figure}[htbp!]
\centering
\includegraphics[width=\textwidth]{logo.png}
\caption[]{\noindent This is \emph{our} logo}
\end{figure}


Remote urls are not permitted. Also, the path checking is rather
strict, so please use just alphanumeric filenames for your images.

\section{Floating images and adjusting width}


By default, image width in the HTML output is kept to the original
one, whatever it is, but limited via CSS, while in the PDF expands to
fill the page width. This creates problems.


Starting with version 0.07, it's possible to set the width of the
image, \textbf{in percent}, just appending it to the file name after a
whitespace.


It's also possible to make it a right or left float, adding the
character \texttt{l} or \texttt{r}, or to mark it as a "fullpage" with \texttt{f}.


This one will have a width to 80\% of the page width:


\begin{alltt}
[[logo.png 80]]

\end{alltt}


This one will have a width to 80\% of the page width and will be a left
float.


\begin{alltt}
[[logo.png 80 l]]

\end{alltt}


This one will have a width to 80\% of the page width and will be a right
float.


\begin{alltt}
[[logo.png 80 r]]

\end{alltt}


The Following will be marked as a fullpage.


\begin{alltt}
[[logo.png f]]

\end{alltt}


Floating without a width doesn't make sense, as in the PDF will be no
room for wrapping text around it.

\subsection{Examples}


The next figures use these codes:


\begin{alltt}
[[logo.png 10r][right float]]
[[logo.png 10l][left float]]
[[logo.png 50f][full page]]

\end{alltt}


Suggestions that I write my memoirs came to me when I had barely begun
to live, and continued all through the years. But I never paid heed to
the proposal. I was living my life intensely — what need to write
about it? Another reason for my reluctance was the conviction I
entertained that one should write about one’s life only when one had
ceased to stand in the very torrent of it. “When one has reached a
good philosophic age,” I used to tell my friends, “capable of viewing
the tragedies and comedies of life impersonally and detachedly —
particularly one’s own life — one is likely to create an autobiography
worth while.” Still feeling adolescently young in spite of advancing
years, I did not consider myself competent to undertake such a task.
Moreover, I always lacked the necessary leisure for concentrated
writing.


\begin{wrapfigure}{r}{0.10\textwidth}
\centering
\includegraphics[width=0.10\textwidth]{logo.png}
\caption[]{\noindent right float}
\end{wrapfigure}


My enforced European inactivity left me enough time to read a great
deal, including biographies and autobiographies. I discovered, much to
my discomfiture, that old age, far from ripening wisdom and


mellowness, is too often fraught with senility, narrowness, and petty
rancour. I would not risk such a calamity, and I began to think
seriously about writing my life.


\begin{wrapfigure}{l}{0.10\textwidth}
\centering
\includegraphics[width=0.10\textwidth]{logo.png}
\caption[]{\noindent left float}
\end{wrapfigure}


The great difficulty that faced me was lack of historical data for my
work. Almost everything in the way of books, correspondence, and
similar material that I had accumulated during the thirty-five years


of my life in the United States had been confiscated by the Department
of Justice raiders and never returned. I lacked even my personal set
of the Mother Earth magazine, which I had published for twelve years.


It was a problem I could see no solution for. Sceptic that I am, I had
overlooked the magic power of friendship, which had so often in my
life made mountains move. My staunch friends Leonard D. Abbott, Agnes
Inglis, W. S. Van Valkenburgh, and others soon put my doubts to shame.


Agnes, the founder of the Labadie Library in Detroit, containing the
richest collection of radical and revolutionary material in America,
came to my aid with her usual readiness. Leonard did his share, and
Van spent all his free time in research work for me.


The great difficulty that faced me was lack of historical data for my
work. Almost everything in the way of books, correspondence, and
similar material that I had accumulated during the thirty-five years


of my life in the United States had been confiscated by the Department
of Justice raiders and never returned. I lacked even my personal set
of the Mother Earth magazine, which I had published for twelve years.


It was a problem I could see no solution for. Sceptic that I am, I had
overlooked the magic power of friendship, which had so often in my
life made mountains move. My staunch friends Leonard D. Abbott, Agnes
Inglis, W. S. Van Valkenburgh, and others soon put my doubts to shame.


Agnes, the founder of the Labadie Library in Detroit, containing the
richest collection of radical and revolutionary material in America,
came to my aid with her usual readiness. Leonard did his share, and
Van spent all his free time in research work for me.


\begin{figure}[htbp!]
\centering
\includegraphics[width=0.50\textwidth]{logo.png}
\caption[]{\noindent full page}
\end{figure}
\clearpage

\chapter{Inserting a horizontal line or anchor}

\section{Horizontal Rules}


Four or more dashes indicate a horizontal rule.  Be sure to put blank
lines around it, or it will be considered part of the proceeding or
following paragraph, like this: ---------.


Example


\begin{alltt}
----

\end{alltt}


Results:


\hairline

\chapter{Lines to omit from published output}


\begin{alltt}
; Comment text goes here.

\end{alltt}



\begin{comment}
; this will disappear completely
\end{comment}


That is, only a semi-colon at the beginning of a line, followed by a
literal space, will cause that line to be treated as a comment and
\textbf{totally removed from the output.}


You can alternatively surround the region with the \texttt{<comment>} tag. In
this case the HTML output will be wrapped in a \texttt{<div>} with the
\texttt{display:none} property, so it can be turned visibile changing the
CSS.


\begin{alltt}
<comment>
This won't be published, but in the HTML is there, only hidden
</comment>

\end{alltt}



\begin{comment}

This won't be published, but in the HTML is there, only hidden


\end{comment}

\chapter{Plays and bibliographies}


Unlike the original Muse, this implementation doesn't support external
sources for citations, but provides an environment which can used to
compose the list of cited works without resorting to lists, wrapping
all in a <biblio> tag.


\begin{alltt}
<biblio>
The author, *Title*, published on
[[http://theanarchistlibrary.org][The Anarchist Library]], with a
very, very, very, very, very, very, very, very, very, very, very,
very, long description

Another author, *Another title*, published in the real world. with a
very, very, very, very, very, very, very, very, very, very, very,
very, long description
</biblio>

\end{alltt}



\begin{amusebiblio}


The author, \emph{Title}, published on
\href{http://theanarchistlibrary.org}{The Anarchist Library}, with a
very, very, very, very, very, very, very, very, very, very, very,
very, long description


Another author, \emph{Another title}, published in the real world. with a
very, very, very, very, very, very, very, very, very, very, very,
very, long description



\end{amusebiblio}


The list is wrapped in a \texttt{<div>} with class set to biblio, while the TeX
source get wrapped in \texttt{\textbackslash{}startawikibiblio} and \texttt{\textbackslash{}stopawikibiblio}. You are
free to change class and definition. A reversed indentation is
suggested.


The same goes with the \texttt{<play>} environment, which is supposed to wrap
theathrical plays, when you want a reverse indentation and more
spacing between the paragraph (without resorting to \texttt{<verse>} and
hardcoding the spaces.


\begin{alltt}
<play>
**Pol.** Ophelia, walke you heere. Gracious so please ye <br>
We will bestow our selues: Reade on this booke, <br>
That shew of such an exercise may colour <br>
Your lonelinesse. We are oft too blame in this, <br>
'Tis too much prou'd, that with Deuotions visage, <br>
And pious Action, we do surge o're <br>
The diuell himselfe

**King.** Oh 'tis true: <br>
How smart a lash that speech doth giue my Conscience? <br>
The Harlots Cheeke beautied with plaist'ring Art <br>
Is not more vgly to the thing that helpes it, <br>
Then is my deede, to my most painted word. <br>
Oh heauie burthen!

</play>

\end{alltt}



\begin{amuseplay}


\textbf{Pol.} Ophelia, walke you heere. Gracious so please ye \forcelinebreak 
We will bestow our selues: Reade on this booke, \forcelinebreak 
That shew of such an exercise may colour \forcelinebreak 
Your lonelinesse. We are oft too blame in this, \forcelinebreak 
'Tis too much prou'd, that with Deuotions visage, \forcelinebreak 
And pious Action, we do surge o're \forcelinebreak 
The diuell himselfe


\textbf{King.} Oh 'tis true: \forcelinebreak 
How smart a lash that speech doth giue my Conscience? \forcelinebreak 
The Harlots Cheeke beautied with plaist'ring Art \forcelinebreak 
Is not more vgly to the thing that helpes it, \forcelinebreak 
Then is my deede, to my most painted word. \forcelinebreak 
Oh heauie burthen!



\end{amuseplay}

\part{GNU Free Documentation License}


\begin{alltt}
                GNU Free Documentation License
                 Version 1.3, 3 November 2008


 Copyright (C) 2000, 2001, 2002, 2007, 2008 Free Software Foundation, Inc.
     <http://fsf.org/>
 Everyone is permitted to copy and distribute verbatim copies
 of this license document, but changing it is not allowed.

0. PREAMBLE

The purpose of this License is to make a manual, textbook, or other
functional and useful document "free" in the sense of freedom: to
assure everyone the effective freedom to copy and redistribute it,
with or without modifying it, either commercially or noncommercially.
Secondarily, this License preserves for the author and publisher a way
to get credit for their work, while not being considered responsible
for modifications made by others.

This License is a kind of "copyleft", which means that derivative
works of the document must themselves be free in the same sense.  It
complements the GNU General Public License, which is a copyleft
license designed for free software.

We have designed this License in order to use it for manuals for free
software, because free software needs free documentation: a free
program should come with manuals providing the same freedoms that the
software does.  But this License is not limited to software manuals;
it can be used for any textual work, regardless of subject matter or
whether it is published as a printed book.  We recommend this License
principally for works whose purpose is instruction or reference.


1. APPLICABILITY AND DEFINITIONS

This License applies to any manual or other work, in any medium, that
contains a notice placed by the copyright holder saying it can be
distributed under the terms of this License.  Such a notice grants a
world-wide, royalty-free license, unlimited in duration, to use that
work under the conditions stated herein.  The "Document", below,
refers to any such manual or work.  Any member of the public is a
licensee, and is addressed as "you".  You accept the license if you
copy, modify or distribute the work in a way requiring permission
under copyright law.

A "Modified Version" of the Document means any work containing the
Document or a portion of it, either copied verbatim, or with
modifications and/or translated into another language.

A "Secondary Section" is a named appendix or a front-matter section of
the Document that deals exclusively with the relationship of the
publishers or authors of the Document to the Document's overall
subject (or to related matters) and contains nothing that could fall
directly within that overall subject.  (Thus, if the Document is in
part a textbook of mathematics, a Secondary Section may not explain
any mathematics.)  The relationship could be a matter of historical
connection with the subject or with related matters, or of legal,
commercial, philosophical, ethical or political position regarding
them.

The "Invariant Sections" are certain Secondary Sections whose titles
are designated, as being those of Invariant Sections, in the notice
that says that the Document is released under this License.  If a
section does not fit the above definition of Secondary then it is not
allowed to be designated as Invariant.  The Document may contain zero
Invariant Sections.  If the Document does not identify any Invariant
Sections then there are none.

The "Cover Texts" are certain short passages of text that are listed,
as Front-Cover Texts or Back-Cover Texts, in the notice that says that
the Document is released under this License.  A Front-Cover Text may
be at most 5 words, and a Back-Cover Text may be at most 25 words.

A "Transparent" copy of the Document means a machine-readable copy,
represented in a format whose specification is available to the
general public, that is suitable for revising the document
straightforwardly with generic text editors or (for images composed of
pixels) generic paint programs or (for drawings) some widely available
drawing editor, and that is suitable for input to text formatters or
for automatic translation to a variety of formats suitable for input
to text formatters.  A copy made in an otherwise Transparent file
format whose markup, or absence of markup, has been arranged to thwart
or discourage subsequent modification by readers is not Transparent.
An image format is not Transparent if used for any substantial amount
of text.  A copy that is not "Transparent" is called "Opaque".

Examples of suitable formats for Transparent copies include plain
ASCII without markup, Texinfo input format, LaTeX input format, SGML
or XML using a publicly available DTD, and standard-conforming simple
HTML, PostScript or PDF designed for human modification.  Examples of
transparent image formats include PNG, XCF and JPG.  Opaque formats
include proprietary formats that can be read and edited only by
proprietary word processors, SGML or XML for which the DTD and/or
processing tools are not generally available, and the
machine-generated HTML, PostScript or PDF produced by some word
processors for output purposes only.

The "Title Page" means, for a printed book, the title page itself,
plus such following pages as are needed to hold, legibly, the material
this License requires to appear in the title page.  For works in
formats which do not have any title page as such, "Title Page" means
the text near the most prominent appearance of the work's title,
preceding the beginning of the body of the text.

The "publisher" means any person or entity that distributes copies of
the Document to the public.

A section "Entitled XYZ" means a named subunit of the Document whose
title either is precisely XYZ or contains XYZ in parentheses following
text that translates XYZ in another language.  (Here XYZ stands for a
specific section name mentioned below, such as "Acknowledgements",
"Dedications", "Endorsements", or "History".)  To "Preserve the Title"
of such a section when you modify the Document means that it remains a
section "Entitled XYZ" according to this definition.

The Document may include Warranty Disclaimers next to the notice which
states that this License applies to the Document.  These Warranty
Disclaimers are considered to be included by reference in this
License, but only as regards disclaiming warranties: any other
implication that these Warranty Disclaimers may have is void and has
no effect on the meaning of this License.

2. VERBATIM COPYING

You may copy and distribute the Document in any medium, either
commercially or noncommercially, provided that this License, the
copyright notices, and the license notice saying this License applies
to the Document are reproduced in all copies, and that you add no
other conditions whatsoever to those of this License.  You may not use
technical measures to obstruct or control the reading or further
copying of the copies you make or distribute.  However, you may accept
compensation in exchange for copies.  If you distribute a large enough
number of copies you must also follow the conditions in section 3.

You may also lend copies, under the same conditions stated above, and
you may publicly display copies.


3. COPYING IN QUANTITY

If you publish printed copies (or copies in media that commonly have
printed covers) of the Document, numbering more than 100, and the
Document's license notice requires Cover Texts, you must enclose the
copies in covers that carry, clearly and legibly, all these Cover
Texts: Front-Cover Texts on the front cover, and Back-Cover Texts on
the back cover.  Both covers must also clearly and legibly identify
you as the publisher of these copies.  The front cover must present
the full title with all words of the title equally prominent and
visible.  You may add other material on the covers in addition.
Copying with changes limited to the covers, as long as they preserve
the title of the Document and satisfy these conditions, can be treated
as verbatim copying in other respects.

If the required texts for either cover are too voluminous to fit
legibly, you should put the first ones listed (as many as fit
reasonably) on the actual cover, and continue the rest onto adjacent
pages.

If you publish or distribute Opaque copies of the Document numbering
more than 100, you must either include a machine-readable Transparent
copy along with each Opaque copy, or state in or with each Opaque copy
a computer-network location from which the general network-using
public has access to download using public-standard network protocols
a complete Transparent copy of the Document, free of added material.
If you use the latter option, you must take reasonably prudent steps,
when you begin distribution of Opaque copies in quantity, to ensure
that this Transparent copy will remain thus accessible at the stated
location until at least one year after the last time you distribute an
Opaque copy (directly or through your agents or retailers) of that
edition to the public.

It is requested, but not required, that you contact the authors of the
Document well before redistributing any large number of copies, to
give them a chance to provide you with an updated version of the
Document.


4. MODIFICATIONS

You may copy and distribute a Modified Version of the Document under
the conditions of sections 2 and 3 above, provided that you release
the Modified Version under precisely this License, with the Modified
Version filling the role of the Document, thus licensing distribution
and modification of the Modified Version to whoever possesses a copy
of it.  In addition, you must do these things in the Modified Version:

A. Use in the Title Page (and on the covers, if any) a title distinct
   from that of the Document, and from those of previous versions
   (which should, if there were any, be listed in the History section
   of the Document).  You may use the same title as a previous version
   if the original publisher of that version gives permission.
B. List on the Title Page, as authors, one or more persons or entities
   responsible for authorship of the modifications in the Modified
   Version, together with at least five of the principal authors of the
   Document (all of its principal authors, if it has fewer than five),
   unless they release you from this requirement.
C. State on the Title page the name of the publisher of the
   Modified Version, as the publisher.
D. Preserve all the copyright notices of the Document.
E. Add an appropriate copyright notice for your modifications
   adjacent to the other copyright notices.
F. Include, immediately after the copyright notices, a license notice
   giving the public permission to use the Modified Version under the
   terms of this License, in the form shown in the Addendum below.
G. Preserve in that license notice the full lists of Invariant Sections
   and required Cover Texts given in the Document's license notice.
H. Include an unaltered copy of this License.
I. Preserve the section Entitled "History", Preserve its Title, and add
   to it an item stating at least the title, year, new authors, and
   publisher of the Modified Version as given on the Title Page.  If
   there is no section Entitled "History" in the Document, create one
   stating the title, year, authors, and publisher of the Document as
   given on its Title Page, then add an item describing the Modified
   Version as stated in the previous sentence.
J. Preserve the network location, if any, given in the Document for
   public access to a Transparent copy of the Document, and likewise
   the network locations given in the Document for previous versions
   it was based on.  These may be placed in the "History" section.
   You may omit a network location for a work that was published at
   least four years before the Document itself, or if the original
   publisher of the version it refers to gives permission.
K. For any section Entitled "Acknowledgements" or "Dedications",
   Preserve the Title of the section, and preserve in the section all
   the substance and tone of each of the contributor acknowledgements
   and/or dedications given therein.
L. Preserve all the Invariant Sections of the Document,
   unaltered in their text and in their titles.  Section numbers
   or the equivalent are not considered part of the section titles.
M. Delete any section Entitled "Endorsements".  Such a section
   may not be included in the Modified Version.
N. Do not retitle any existing section to be Entitled "Endorsements"
   or to conflict in title with any Invariant Section.
O. Preserve any Warranty Disclaimers.

If the Modified Version includes new front-matter sections or
appendices that qualify as Secondary Sections and contain no material
copied from the Document, you may at your option designate some or all
of these sections as invariant.  To do this, add their titles to the
list of Invariant Sections in the Modified Version's license notice.
These titles must be distinct from any other section titles.

You may add a section Entitled "Endorsements", provided it contains
nothing but endorsements of your Modified Version by various
parties--for example, statements of peer review or that the text has
been approved by an organization as the authoritative definition of a
standard.

You may add a passage of up to five words as a Front-Cover Text, and a
passage of up to 25 words as a Back-Cover Text, to the end of the list
of Cover Texts in the Modified Version.  Only one passage of
Front-Cover Text and one of Back-Cover Text may be added by (or
through arrangements made by) any one entity.  If the Document already
includes a cover text for the same cover, previously added by you or
by arrangement made by the same entity you are acting on behalf of,
you may not add another; but you may replace the old one, on explicit
permission from the previous publisher that added the old one.

The author(s) and publisher(s) of the Document do not by this License
give permission to use their names for publicity for or to assert or
imply endorsement of any Modified Version.


5. COMBINING DOCUMENTS

You may combine the Document with other documents released under this
License, under the terms defined in section 4 above for modified
versions, provided that you include in the combination all of the
Invariant Sections of all of the original documents, unmodified, and
list them all as Invariant Sections of your combined work in its
license notice, and that you preserve all their Warranty Disclaimers.

The combined work need only contain one copy of this License, and
multiple identical Invariant Sections may be replaced with a single
copy.  If there are multiple Invariant Sections with the same name but
different contents, make the title of each such section unique by
adding at the end of it, in parentheses, the name of the original
author or publisher of that section if known, or else a unique number.
Make the same adjustment to the section titles in the list of
Invariant Sections in the license notice of the combined work.

In the combination, you must combine any sections Entitled "History"
in the various original documents, forming one section Entitled
"History"; likewise combine any sections Entitled "Acknowledgements",
and any sections Entitled "Dedications".  You must delete all sections
Entitled "Endorsements".


6. COLLECTIONS OF DOCUMENTS

You may make a collection consisting of the Document and other
documents released under this License, and replace the individual
copies of this License in the various documents with a single copy
that is included in the collection, provided that you follow the rules
of this License for verbatim copying of each of the documents in all
other respects.

You may extract a single document from such a collection, and
distribute it individually under this License, provided you insert a
copy of this License into the extracted document, and follow this
License in all other respects regarding verbatim copying of that
document.


7. AGGREGATION WITH INDEPENDENT WORKS

A compilation of the Document or its derivatives with other separate
and independent documents or works, in or on a volume of a storage or
distribution medium, is called an "aggregate" if the copyright
resulting from the compilation is not used to limit the legal rights
of the compilation's users beyond what the individual works permit.
When the Document is included in an aggregate, this License does not
apply to the other works in the aggregate which are not themselves
derivative works of the Document.

If the Cover Text requirement of section 3 is applicable to these
copies of the Document, then if the Document is less than one half of
the entire aggregate, the Document's Cover Texts may be placed on
covers that bracket the Document within the aggregate, or the
electronic equivalent of covers if the Document is in electronic form.
Otherwise they must appear on printed covers that bracket the whole
aggregate.


8. TRANSLATION

Translation is considered a kind of modification, so you may
distribute translations of the Document under the terms of section 4.
Replacing Invariant Sections with translations requires special
permission from their copyright holders, but you may include
translations of some or all Invariant Sections in addition to the
original versions of these Invariant Sections.  You may include a
translation of this License, and all the license notices in the
Document, and any Warranty Disclaimers, provided that you also include
the original English version of this License and the original versions
of those notices and disclaimers.  In case of a disagreement between
the translation and the original version of this License or a notice
or disclaimer, the original version will prevail.

If a section in the Document is Entitled "Acknowledgements",
"Dedications", or "History", the requirement (section 4) to Preserve
its Title (section 1) will typically require changing the actual
title.


9. TERMINATION

You may not copy, modify, sublicense, or distribute the Document
except as expressly provided under this License.  Any attempt
otherwise to copy, modify, sublicense, or distribute it is void, and
will automatically terminate your rights under this License.

However, if you cease all violation of this License, then your license
from a particular copyright holder is reinstated (a) provisionally,
unless and until the copyright holder explicitly and finally
terminates your license, and (b) permanently, if the copyright holder
fails to notify you of the violation by some reasonable means prior to
60 days after the cessation.

Moreover, your license from a particular copyright holder is
reinstated permanently if the copyright holder notifies you of the
violation by some reasonable means, this is the first time you have
received notice of violation of this License (for any work) from that
copyright holder, and you cure the violation prior to 30 days after
your receipt of the notice.

Termination of your rights under this section does not terminate the
licenses of parties who have received copies or rights from you under
this License.  If your rights have been terminated and not permanently
reinstated, receipt of a copy of some or all of the same material does
not give you any rights to use it.


10. FUTURE REVISIONS OF THIS LICENSE

The Free Software Foundation may publish new, revised versions of the
GNU Free Documentation License from time to time.  Such new versions
will be similar in spirit to the present version, but may differ in
detail to address new problems or concerns.  See
http://www.gnu.org/copyleft/.

Each version of the License is given a distinguishing version number.
If the Document specifies that a particular numbered version of this
License "or any later version" applies to it, you have the option of
following the terms and conditions either of that specified version or
of any later version that has been published (not as a draft) by the
Free Software Foundation.  If the Document does not specify a version
number of this License, you may choose any version ever published (not
as a draft) by the Free Software Foundation.  If the Document
specifies that a proxy can decide which future versions of this
License can be used, that proxy's public statement of acceptance of a
version permanently authorizes you to choose that version for the
Document.

11. RELICENSING

"Massive Multiauthor Collaboration Site" (or "MMC Site") means any
World Wide Web server that publishes copyrightable works and also
provides prominent facilities for anybody to edit those works.  A
public wiki that anybody can edit is an example of such a server.  A
"Massive Multiauthor Collaboration" (or "MMC") contained in the site
means any set of copyrightable works thus published on the MMC site.

"CC-BY-SA" means the Creative Commons Attribution-Share Alike 3.0
license published by Creative Commons Corporation, a not-for-profit
corporation with a principal place of business in San Francisco,
California, as well as future copyleft versions of that license
published by that same organization.

"Incorporate" means to publish or republish a Document, in whole or in
part, as part of another Document.

An MMC is "eligible for relicensing" if it is licensed under this
License, and if all works that were first published under this License
somewhere other than this MMC, and subsequently incorporated in whole or
in part into the MMC, (1) had no cover texts or invariant sections, and
(2) were thus incorporated prior to November 1, 2008.

The operator of an MMC Site may republish an MMC contained in the site
under CC-BY-SA on the same site at any time before August 1, 2009,
provided the MMC is eligible for relicensing.


ADDENDUM: How to use this License for your documents

To use this License in a document you have written, include a copy of
the License in the document and put the following copyright and
license notices just after the title page:

    Copyright (c)  YEAR  YOUR NAME.
    Permission is granted to copy, distribute and/or modify this document
    under the terms of the GNU Free Documentation License, Version 1.3
    or any later version published by the Free Software Foundation;
    with no Invariant Sections, no Front-Cover Texts, and no Back-Cover Texts.
    A copy of the license is included in the section entitled "GNU
    Free Documentation License".

If you have Invariant Sections, Front-Cover Texts and Back-Cover Texts,
replace the "with...Texts." line with this:

    with the Invariant Sections being LIST THEIR TITLES, with the
    Front-Cover Texts being LIST, and with the Back-Cover Texts being LIST.

If you have Invariant Sections without Cover Texts, or some other
combination of the three, merge those two alternatives to suit the
situation.

If your document contains nontrivial examples of program code, we
recommend releasing these examples in parallel under your choice of
free software license, such as the GNU General Public License,
to permit their use in free software.


\end{alltt}



\clearpage
% new page for the colophon

\thispagestyle{empty}

\begin{center}





\strut

\end{center}


\strut


\vfill

\begin{center}



The Text::Amuse markup manual

The writer's guide

2012

\bigskip







% Here an URL maybe?

\end{center}

\end{document}

